\section{任务定义与记号}
\label{sec:problem}

\paragraph{输入。}核心输入是一份 A 股上市公司年报（PDF），可附季报、业绩预告、卖方一致预期、可比公司列表和一份待核查的研报草稿。对研报生成任务，还允许系统联网检索年报之后的公开信息，但检索结果必须登记（第 \ref{sec:sources} 节）。

\paragraph{文档与片段。}解析器把文档 $D$ 切成带锚点的片段序列 $\{c_k\}$，每个片段记录文本、表格行、所在页码、表格标题、单位提示，以及（若为跨页续表）它所接续的片段。片段是来源定位的最小单位。

\paragraph{事实与真值表。}设指标集合为 $\mathcal{M}$（当前 49 个，涵盖利润表、资产负债表、现金流量表核心科目与主要比率），期间集合 $\mathcal{P}$ 包括年度 \code{2025}、半年度 \code{2025H1}、单季 \code{2025Q3} 与累计 \code{2025Q1-3}。一条\textbf{事实}是六元组
\begin{equation}
  f = (m,\; p,\; v,\; r,\; u,\; \sigma),\qquad m\in\mathcal{M},\; p\in\mathcal{P},
\end{equation}
其中 $v$ 是换算到基本单位后的数值（金额为元，比率为百分数值，如 $14.49$ 表示 $14.49\%$），$r$ 是原文写法（如 \code{"1,234.50"}），$u$ 是原文单位，$\sigma=(\text{文件},\text{页码},\text{片段},\text{行},\text{列})$ 是来源定位。程序派生的值（例如单季拆分）不带 $\sigma$，而带公式与所依赖的事实。所有事实与派生值构成\textbf{真值表} $T:\mathcal{M}\times\mathcal{P}\rightharpoonup f$。

\paragraph{任务 A：财报分析（赛题 2）。}给定 $T$，输出同比、环比、单季拆分、勾稽校验、异常信号、利润变动归因、实际与一致预期或业绩预告的对比，以及每股收益质量判断。所有数值都是 $T$ 的确定性函数。

\paragraph{任务 B：估值建模（赛题 4）。}给定 $T$、经营数据 $E$（分产品量价、产能等，逐格核对过原文）与估值假设 $\Theta$（预测参数、折现参数、参考股价及其来源），输出三情景三表模型、FCFF 折现估值、敏感性矩阵、可比公司估值、反向折现（股价隐含了什么）和程序评级。$\Theta$ 中每一项都标明为“假设”，并写出依据。

\paragraph{任务 C：研报纠错（赛题 5）。}给定研报草稿 $R$，抽取可核验的数字断言 $a=(m,p,\tau,x,r,u)$，其中 $\tau\in\{\text{水平值},\text{比率},\text{同比},\text{百分点变动}\}$，并把每条断言判为
\[
  \{\text{正确},\ \text{数值错误},\ \text{单位错误},\ \text{期间错配},\ \text{口径错误},\ \text{方向错误},\ \text{无法核实}\},
\]
同时给出真值、来源定位与修改建议。

\paragraph{研报生成。}在完成 A、B 之后，生成一份卖方体例的深度研报文稿 $N$ 与排版后的 PDF。记 $\mathrm{Num}(N)$ 为 $N$ 中可识别的“数字 + 单位”记号集合，$V$ 为本次运行结论性产物中的数值集合，$S$ 为外部来源登记表中的数值集合。我们要求
\begin{equation}
  \forall\, x\in \mathrm{Num}(N):\quad \exists\, v\in V\cup S,\quad |x - \phi_u(v)| \le \delta(x),
  \label{eq:provenance}
\end{equation}
其中 $\phi_u$ 把基本单位换算到 $x$ 的展示单位，$\delta(x)$ 是由 $x$ 的展示精度决定的容差（第 \ref{sec:guard} 节）；若 $v$ 只来自 $S$，则 $x$ 所在的句子必须带对应的来源编号。式 \eqref{eq:provenance} 是硬约束；研报的分析质量与表达质量是软目标，由第 \ref{sec:eval} 节的评分标准度量。

\paragraph{单位约定。}财务事实内部统一用元；估值模块（三表、DCF、可比公司）用百万元，每股用元；比率存百分数值。同一页出现多个“单位：”声明时，每张表取其上方最近的声明。这些约定在代码与技能说明中保持一致，护栏在比较时按约定换算。
